\documentclass[12pt,a4paper]{report}

% --- Language & Font Support ---
\usepackage{polyglossia}
\usepackage{subcaption}
\setmainlanguage{english}
\setotherlanguage{arabic}

\newfontfamily\arabicfont[Script=Arabic]{Amiri}

% --- Base Packages ---
\usepackage{geometry}
\geometry{margin=2.5cm}
\usepackage{graphicx}
\usepackage{hyperref}
\usepackage{xcolor}
\usepackage{titlesec}
\usepackage{tabularx}
\usepackage{booktabs}
\usepackage{lstlisting}
\usepackage{listings}
\usepackage{fancyhdr}
\usepackage{float}

% --- Code Block Configuration ---
\definecolor{codeblue}{rgb}{0.13,0.13,1}
\definecolor{codegreen}{rgb}{0,0.6,0}
\definecolor{codegray}{rgb}{0.5,0.5,0.5}
\definecolor{backcolour}{rgb}{0.97,0.97,0.97}

\lstset{
    backgroundcolor=\color{backcolour},   
    commentstyle=\color{codegreen},
    keywordstyle=\color{codeblue},
    numberstyle=\tiny\color{codegray},
    stringstyle=\color{orange},
    basicstyle=\ttfamily\footnotesize,
    breaklines=true,                 
    captionpos=b,                    
    keepspaces=true,                 
    numbers=left,                    
    tabsize=2,
    showstringspaces=false
}

% --- Header and Footer ---
\pagestyle{fancy}
\fancyhf{}
\fancyhead[L]{Frontend Architecture \& Core Template - Group 8}
\fancyfoot[C]{\thepage}

\begin{document}

% --- Cover Page ---
\begin{titlepage}
    \centering
    
    {\Large \textbf{\textarabic{الجمهورية الجزائرية الديمقراطية الشعبية}}} \\
    {\small \textbf{\textarabic{وزارة التعليم العالي والبحث العلمي}}} \\[0.3cm]
    
    \Large \textbf{PEOPLE'S DEMOCRATIC REPUBLIC OF ALGERIA}\\
    \small \textbf{MINISTRY OF HIGHER EDUCATION AND SCIENTIFIC RESEARCH}
    
    \vfill

    \includegraphics[width=0.25\textwidth]{logo_univ.png} % Suggestion: Place logo here
    \\[0.4cm]
    
    {\Large \textbf{\textarabic{جامعة ابن خلدون تيارت}}} \\
    \Large \textbf{Ibn Khaldoun University of Tiaret}\\
    \large Department of Computer Science
    
    \vfill
    
    \huge \textbf{Final Year Project (L3)}\\[0.5cm]
    \LARGE \textbf{Development of a Frontend Architecture and a Modular Design System}
    
    \vfill
    
    \large \textbf{Presented by:}\\[0.5cm]
    \Large 
    \begin{tabular}{c}
        \textbf{Djellil Abdelkader Charef Eddine} \\
        \textbf{Gaid Oussama} \\
        \textbf{Derkaoui Mohamed} \\
    \end{tabular}
    
    \vfill
    
    \large \textbf{Group 8} \\
    (Frontend Architecture \& Core Template)
    
    \vfill
    
    \large \textbf{Academic Year:} 2025 - 2026
\end{titlepage}

\selectlanguage{english}

\chapter{Module Design: Core Template \& Frontend Architecture}

% ---------------------------------------------------------
% 3.1 Presentation, Objectives, Features, and Constraints
% ---------------------------------------------------------
\section{Module Presentation, Objectives, Features, Constraints, and Communication}

\subsection{Presentation and Objectives}
The \textbf{Core Template} module constitutes the foundational infrastructure of the university platform. It is strictly dedicated to standardizing the user interface (UI), user experience (UX), and providing a global configuration shell via the \texttt{SiteSetting} model. The primary objective is to deliver a robust \textit{Design System} (a collection of reusable UI components) and a standard dashboard layout, enabling other functional groups to build their specific logic without reinventing basic interface elements.

\subsection{Features and Constraints}
\begin{itemize}
    \item \textbf{Theming Engine:} Dynamic management of Light/Dark modes and configurable institutional accent colors using semantic CSS variables.
    \item \textbf{Comprehensive UI Library:} Provision of reusable atomic components such as \texttt{Button}, \texttt{Input}, \texttt{Alert}, \texttt{Modal}, and \texttt{KpiCard}.
    \item \textbf{Global Settings Management:} A dedicated backend model (\texttt{SiteSetting}) allowing administrators to dynamically alter the platform's visual identity (logo, title, theme).
    \item \textbf{Technical Constraints:} High performance rendering through optimized React components and native bidirectionality support (LTR/RTL) catering to bilingual requirements.
\end{itemize}

\subsection{Communication (API Contract)}
A core contribution of Group 8 is standardizing how the client application communicates with the backend. We implemented an Axios interceptor layer that guarantees every API response—whether successful or faulty—follows a strict JSON contract. This uniformity simplifies error handling for all developer teams across the platform.

An example of our standardized JSON response contract:
\begin{lstlisting}[language=json, caption=Standard API Response JSON Contract]
{
  "success": true,
  "data": {
    "siteTitle": "University Portal",
    "theme": "light",
    "logoUrl": "/uploads/logo.png"
  },
  "message": "Settings retrieved successfully"
}
\end{lstlisting}

\begin{figure}[H]
    \centering
    % TODO: Add a screenshot of Postman or browser console showing a successful API response here
    \includegraphics[width=0.8\textwidth]{screenshots/api-response-postman.png}
    \caption{API response interception and standardization in action.}
    \label{fig:api_contract}
\end{figure}

% ---------------------------------------------------------
% 3.2 Modeling and Design
% ---------------------------------------------------------
\section{Modeling and Design}

\subsection{Global Configuration Entity (\texttt{SiteSetting})}
While other groups managed complex business domains (such as assignments or users), Group 8 focused strictly on the global configuration layer. We developed the \texttt{SiteSetting} backend module, which acts as a dynamic singleton. This entity controls global properties accessed by the entire frontend upon loading.

\begin{figure}[H]
    \centering
    % TODO: Add a screenshot of the SiteSettings Admin Page UI
    \includegraphics[width=0.8\textwidth]{screenshots/siteconfig-admin.png} 
    \caption{Administration Interface: Managing the SiteSetting database model dynamically.}
    \label{fig:admin_siteconfig}
\end{figure}


% ---------------------------------------------------------
% 3.3 Implementation and Integration
% ---------------------------------------------------------
\section{Implementation and Integration}

\subsection{Design Tokens and the \texttt{ThemeProvider}}
The UI relies heavily on semantic Design Tokens via Tailwind CSS. Instead of hardcoding colors, the styles map to generic variables like \texttt{--color-surface} and \texttt{--color-ink}. Our \texttt{ThemeProvider} component manages these variables within React's Context, directly manipulating the DOM's root element attributes to facilitate real-time, global theme switching.

\subsection{Atomic Design System Library}
The core deliverable for the frontend is the extensive component library located under \texttt{frontend/src/design-system/components/}. Every component is built to be stateless where possible, relying entirely on provided properties (\texttt{props}) for variants and states.

\begin{figure}[H]
    \centering
    % TODO: Add a screenshot of the "Component Showcase" page showing Buttons/Inputs/Alerts in action
    \includegraphics[width=0.9\textwidth]{screenshots/design-system-showcase.png}
    \caption{Component Showcase: A centralized library for reusable elements.}
    \label{fig:component_library}
\end{figure}

\subsection{The Dashboard Shell (Structural Canvas)}
To unify the user experience across all modules, Group 8 developed the global Dashboard layout. This architectural shell includes:
\begin{itemize}
    \item \textbf{Navigation (\texttt{Sidebar.jsx}):} A collapsible lateral menu providing route structure.
    \item \textbf{Context Bar (\texttt{Topbar.jsx}):} A superior bar containing global actions (Theme toggle, Profile access).
    \item \textbf{Primitive Widgets (\texttt{KpiCard.jsx}):} Reusable analytical blocks (e.g., statistics tiles).
\end{itemize}
\textit{Note: Group 8 provides the structural shell and components. The specific dashboards (Teacher, Admin, Student) and their business logic were populated by respective functional groups.}

\begin{figure}[H]
    \centering
    % TODO: Add a wireframe-style screenshot or a clean view of the empty dashboard shell featuring Sidebar & Topbar
    \includegraphics[width=0.9\textwidth]{screenshots/dashboard-shell.png} 
    \caption{The visual dashboard shell serving as a canvas for other modules.}
    \label{fig:nav_layout}
\end{figure}

% ---------------------------------------------------------
% 3.4 Tests and Validation
% ---------------------------------------------------------
\section{Tests and Validation}

\subsection{Component Isolation and Robustness}
To ensure the resilience of the Design System, components were validated in isolation through an integrated showcase environment. This environment verifies visual regressions and state behavior without requiring a full backend lifecycle.

\subsection{Responsiveness and Adaptability (Responsive Design)}
A major technical constraint managed by Group 8 was the execution of a multi-device UI. Utilizing Tailwind's breakpoint system, components seamlessly jump between column layouts on desktop and stacked views on mobile. Specifically, the \textbf{Sidebar} collapses intelligently into a drawer, preserving screen real estate on mobile devices.

\begin{figure}[H]
    \centering
    % TODO: Add two side-by-side screenshots showing a page on Desktop vs Mobile
    \begin{subfigure}[b]{0.5\textwidth}
        \centering
        \includegraphics[width=\textwidth]{screenshots/responsive-desktop.png}
        \caption{Desktop View}
    \end{subfigure}
    \hspace{0.5cm} 
    \begin{subfigure}[b]{0.25\textwidth}
        \centering
        \includegraphics[width=\textwidth]{screenshots/responsive-mobile.png}
        \caption{Mobile View}
    \end{subfigure}
    
    \caption{Responsive behavior executed by the Core Template grids.}
    \label{fig:responsive_validation}
\end{figure}


% ---------------------------------------------------------
% 3.5 Results
% ---------------------------------------------------------
\section{Results}
Implementing this central architecture delivered significant operational improvements across the development lifecycle:
\begin{itemize}
    \item \textbf{Eliminated Redundancy:} The component wrapper system completely removed the need for other groups to write custom CSS or manage complex UI states (e.g., hover effects, focus rings).
    \item \textbf{Guaranteed Consistency:} Due to the \texttt{ThemeProvider} and normalized tokens, the transition from Light to Dark mode requires zero extra code from functional module developers.
    \item \textbf{Reliable Communication:} By standardizing the Axios layer, payload discrepancies between the client interface and modular backend controllers were effectively eliminated.
\end{itemize}

% ---------------------------------------------------------
% 3.6 Conclusion for Group 8
% ---------------------------------------------------------
\section{Conclusion for Group 8}
Group 8 successfully delivered the architectural skeleton necessary and sufficient for the university's full platform. By abstracting away the complexities of state-driven UI (Design System), global theming, base layout organization (Dashboards), and API communication (Axios interceptors & \texttt{SiteSetting}), we provided the other teams highly-structured terrain built over a robust containerized infrastructure. This modular approach accelerated overall team velocity and ensured a highly professional and robust final interface.

\end{document}